\documentclass{article}
%DIF LATEXDIFF DIFFERENCE FILE
%DIF DEL build/theissen_0919/main_before.tex   Sat Sep 19 16:46:43 2026
%DIF ADD main.tex                              Sat Sep 19 17:37:20 2026

\usepackage{neurips_2026}

\usepackage[utf8]{inputenc}
\usepackage[T1]{fontenc}
\usepackage{hyperref}
\usepackage{url}
\usepackage{booktabs}
\usepackage{amsfonts}
\usepackage{amsmath}
\usepackage{amssymb}   % \lesssim in Appendix A; amsmath alone does not define it
\usepackage{nicefrac}
\usepackage{microtype}
\usepackage{xcolor}
\usepackage{graphicx}
\usepackage{float}     % [H] for the Appendix C figure, so it stays with its paragraph
\setlength{\abovecaptionskip}{4pt}   % float spacing tightened to hold the body on four pages
\setlength{\textfloatsep}{10pt}
\setlength{\floatsep}{8pt}

% The workshop prescribes the first-page footer text. The style file is otherwise untouched.
\makeatletter
\renewcommand{\@noticestring}{%
  Submitted to the 9th Workshop on Machine Learning and the Physical Sciences (ML4PS 2026).
  Do not distribute.%
}
\makeatother

\title{What does self-supervision add to engineered features\\for stellar light curves?}

\author{Anonymous Author(s)\\Affiliation\\Address\\\texttt{email}}

 %DIF > 
%DIF PREAMBLE EXTENSION ADDED BY LATEXDIFF
%DIF UNDERLINE PREAMBLE %DIF PREAMBLE
\RequirePackage[normalem]{ulem} %DIF PREAMBLE
\RequirePackage{color}\definecolor{RED}{rgb}{1,0,0}\definecolor{BLUE}{rgb}{0,0,1} %DIF PREAMBLE
\providecommand{\DIFaddtex}[1]{{\protect\color{blue}\uwave{#1}}} %DIF PREAMBLE
\providecommand{\DIFdeltex}[1]{{\protect\color{red}\sout{#1}}} %DIF PREAMBLE
%DIF SAFE PREAMBLE %DIF PREAMBLE
\providecommand{\DIFaddbegin}{} %DIF PREAMBLE
\providecommand{\DIFaddend}{} %DIF PREAMBLE
\providecommand{\DIFdelbegin}{} %DIF PREAMBLE
\providecommand{\DIFdelend}{} %DIF PREAMBLE
\providecommand{\DIFmodbegin}{} %DIF PREAMBLE
\providecommand{\DIFmodend}{} %DIF PREAMBLE
%DIF FLOATSAFE PREAMBLE %DIF PREAMBLE
\providecommand{\DIFaddFL}[1]{\DIFadd{#1}} %DIF PREAMBLE
\providecommand{\DIFdelFL}[1]{\DIFdel{#1}} %DIF PREAMBLE
\providecommand{\DIFaddbeginFL}{} %DIF PREAMBLE
\providecommand{\DIFaddendFL}{} %DIF PREAMBLE
\providecommand{\DIFdelbeginFL}{} %DIF PREAMBLE
\providecommand{\DIFdelendFL}{} %DIF PREAMBLE
%DIF HYPERREF PREAMBLE %DIF PREAMBLE
\providecommand{\DIFadd}[1]{\texorpdfstring{\DIFaddtex{#1}}{#1}} %DIF PREAMBLE
\providecommand{\DIFdel}[1]{\texorpdfstring{\DIFdeltex{#1}}{}} %DIF PREAMBLE
\providecommand{\DIFscaledelfig}{0.5}
%DIF HIGHLIGHTGRAPHICS PREAMBLE %DIF PREAMBLE
\RequirePackage{settobox} %DIF PREAMBLE
\RequirePackage{letltxmacro} %DIF PREAMBLE
\newsavebox{\DIFdelgraphicsbox} %DIF PREAMBLE
\newlength{\DIFdelgraphicswidth} %DIF PREAMBLE
\newlength{\DIFdelgraphicsheight} %DIF PREAMBLE
% store original definition of \includegraphics %DIF PREAMBLE
\LetLtxMacro{\DIFOincludegraphics}{\includegraphics} %DIF PREAMBLE
\providecommand{\DIFaddincludegraphics}[2][]{{\color{blue}\fbox{\DIFOincludegraphics[#1]{#2}}}} %DIF PREAMBLE
\providecommand{\DIFdelincludegraphics}[2][]{% %DIF PREAMBLE
\sbox{\DIFdelgraphicsbox}{\DIFOincludegraphics[#1]{#2}}% %DIF PREAMBLE
\settoboxwidth{\DIFdelgraphicswidth}{\DIFdelgraphicsbox} %DIF PREAMBLE
\settoboxtotalheight{\DIFdelgraphicsheight}{\DIFdelgraphicsbox} %DIF PREAMBLE
\scalebox{\DIFscaledelfig}{% %DIF PREAMBLE
\parbox[b]{\DIFdelgraphicswidth}{\usebox{\DIFdelgraphicsbox}\\[-\baselineskip] \rule{\DIFdelgraphicswidth}{0em}}\llap{\resizebox{\DIFdelgraphicswidth}{\DIFdelgraphicsheight}{% %DIF PREAMBLE
\setlength{\unitlength}{\DIFdelgraphicswidth}% %DIF PREAMBLE
\begin{picture}(1,1)% %DIF PREAMBLE
\thicklines\linethickness{2pt} %DIF PREAMBLE
{\color[rgb]{1,0,0}\put(0,0){\framebox(1,1){}}}% %DIF PREAMBLE
{\color[rgb]{1,0,0}\put(0,0){\line( 1,1){1}}}% %DIF PREAMBLE
{\color[rgb]{1,0,0}\put(0,1){\line(1,-1){1}}}% %DIF PREAMBLE
\end{picture}% %DIF PREAMBLE
}\hspace*{3pt}}} %DIF PREAMBLE
} %DIF PREAMBLE
\LetLtxMacro{\DIFOaddbegin}{\DIFaddbegin} %DIF PREAMBLE
\LetLtxMacro{\DIFOaddend}{\DIFaddend} %DIF PREAMBLE
\LetLtxMacro{\DIFOdelbegin}{\DIFdelbegin} %DIF PREAMBLE
\LetLtxMacro{\DIFOdelend}{\DIFdelend} %DIF PREAMBLE
\DeclareRobustCommand{\DIFaddbegin}{\DIFOaddbegin \let\includegraphics\DIFaddincludegraphics} %DIF PREAMBLE
\DeclareRobustCommand{\DIFaddend}{\DIFOaddend \let\includegraphics\DIFOincludegraphics} %DIF PREAMBLE
\DeclareRobustCommand{\DIFdelbegin}{\DIFOdelbegin \let\includegraphics\DIFdelincludegraphics} %DIF PREAMBLE
\DeclareRobustCommand{\DIFdelend}{\DIFOaddend \let\includegraphics\DIFOincludegraphics} %DIF PREAMBLE
\LetLtxMacro{\DIFOaddbeginFL}{\DIFaddbeginFL} %DIF PREAMBLE
\LetLtxMacro{\DIFOaddendFL}{\DIFaddendFL} %DIF PREAMBLE
\LetLtxMacro{\DIFOdelbeginFL}{\DIFdelbeginFL} %DIF PREAMBLE
\LetLtxMacro{\DIFOdelendFL}{\DIFdelendFL} %DIF PREAMBLE
\DeclareRobustCommand{\DIFaddbeginFL}{\DIFOaddbeginFL \let\includegraphics\DIFaddincludegraphics} %DIF PREAMBLE
\DeclareRobustCommand{\DIFaddendFL}{\DIFOaddendFL \let\includegraphics\DIFOincludegraphics} %DIF PREAMBLE
\DeclareRobustCommand{\DIFdelbeginFL}{\DIFOdelbeginFL \let\includegraphics\DIFdelincludegraphics} %DIF PREAMBLE
\DeclareRobustCommand{\DIFdelendFL}{\DIFOaddendFL \let\includegraphics\DIFOincludegraphics} %DIF PREAMBLE
%DIF AMSMATHULEM PREAMBLE %DIF PREAMBLE
\makeatletter %DIF PREAMBLE
\let\sout@orig\sout %DIF PREAMBLE
\renewcommand{\sout}[1]{\ifmmode\text{\sout@orig{\ensuremath{#1}}}\else\sout@orig{#1}\fi} %DIF PREAMBLE
\makeatother %DIF PREAMBLE
%DIF COLORLISTINGS PREAMBLE %DIF PREAMBLE
\RequirePackage{listings} %DIF PREAMBLE
\RequirePackage{color} %DIF PREAMBLE
\lstdefinelanguage{DIFcode}{ %DIF PREAMBLE
%DIF DIFCODE_UNDERLINE %DIF PREAMBLE
  moredelim=[il][\color{red}\sout]{\%DIF\ <\ }, %DIF PREAMBLE
  moredelim=[il][\color{blue}\uwave]{\%DIF\ >\ } %DIF PREAMBLE
} %DIF PREAMBLE
\lstdefinestyle{DIFverbatimstyle}{ %DIF PREAMBLE
	language=DIFcode, %DIF PREAMBLE
	basicstyle=\ttfamily, %DIF PREAMBLE
	columns=fullflexible, %DIF PREAMBLE
	keepspaces=true %DIF PREAMBLE
} %DIF PREAMBLE
\lstnewenvironment{DIFverbatim}[1][]{\lstset{style=DIFverbatimstyle}}{} %DIF PREAMBLE
\lstnewenvironment{DIFverbatim*}[1][]{\lstset{style=DIFverbatimstyle,showspaces=true}}{} %DIF PREAMBLE
\lstset{extendedchars=\true,inputencoding=utf8}

%DIF END PREAMBLE EXTENSION ADDED BY LATEXDIFF

\begin{document}

\maketitle

\begin{abstract}
TESS light curves encode a star's pulsation, rotation, binarity and oscillations, \DIFaddbegin \DIFadd{among other physical
properties, }\DIFaddend but the survey
produces far more light curves than can be labelled by hand. Astronomers have engineered a few
dozen physically motivated features from \DIFdelbegin \DIFdel{these light curves}\DIFdelend \DIFaddbegin \DIFadd{them}\DIFaddend , which together remain a strong baseline. Previous work has trained
encoders on light curves and judged them by their own downstream scores, but has not asked what
they add to the engineered features. We pre-train a convolutional encoder on unlabelled TESS light curves to predict each
window's latent representation from those of its neighbours. We then freeze it, add its features to
the 25 engineered features, and read the combination out with one fixed linear probe on ten tasks\DIFdelbegin \DIFdel{spanning variability , asteroseismology and rotation }\DIFdelend \DIFaddbegin \DIFadd{:
five variability classes, four asteroseismic targets and rotation period}\DIFaddend . Adding the learned features
improves the readout on 7 of 10 tasks over the engineered features alone, while an untrained encoder
does not. The combination also outperforms two networks trained end-to-end on the same labelled
stars on most tasks. An encoder trained without labels therefore gives astronomers new features
that complement the ones they already use rather than replace them.
\end{abstract}

\section{Introduction}

A star's light curve carries its pulsation, rotation, eclipses and oscillations, and with them its
age, interior and companions. Wide-field \DIFdelbegin \DIFdel{photometric }\DIFdelend surveys now produce light curves far faster
than they can be labelled. TESS
\citep{ricker2015} \DIFdelbegin \DIFdel{trades sampling rate against sky coverage: }\DIFdelend \DIFaddbegin \DIFadd{returns }\DIFaddend 2-minute \DIFdelbegin \DIFdel{(}\DIFdelend or 20-second \DIFdelbegin \DIFdel{) }\DIFdelend \DIFaddbegin \DIFadd{\mbox{%DIFAUXCMD
\citep{huber2022} }\hskip0pt%DIFAUXCMD
}\DIFaddend light curves for hundreds of
thousands of pre-selected stars, and full-frame images \DIFdelbegin \DIFdel{, every 30 or 10 minutes, for
every }\DIFdelend \DIFaddbegin \DIFadd{of every }\DIFaddend star in the field \DIFdelbegin \DIFdel{\mbox{%DIFAUXCMD
\citep{huber2022}}\hskip0pt%DIFAUXCMD
}\DIFdelend \DIFaddbegin \DIFadd{at cadences that
change between mission cycles (30 minutes, then 10, now 200 seconds)}\DIFaddend . LSST \citep{ivezic2019}
will \DIFdelbegin \DIFdel{widen the gap}\DIFdelend \DIFaddbegin \DIFadd{add }\emph{\DIFadd{billions}} \DIFadd{more light curves, albeit at a cadence of days}\DIFaddend . Self-supervised learning is the
natural response: pre-train one encoder on the unlabelled archive, then attach a cheap readout for
each \DIFdelbegin \DIFdel{downstream }\DIFdelend task. Encoders of this shape exist for ground-based, Gaia and Kepler light curves
\citep{donoso2025astromer2,zuo2025falco,rui2026jepa} and are evaluated by classification accuracy,
sometimes against engineered features \citep{rui2026jepa}; a label-free TESS representation
\citep{poznanski2026} is evaluated by whether repeat observations of one star lie close in its embedding space.
None measures what they add to engineered features.


\clearpage
\bibliographystyle{plainnat}
\bibliography{refs}
\end{document}
